\label{sec:intro}

Route generation is a primitive for mobility simulation, transportation planning, and urban computing~\citep{kong2023mobility_survey,zhang2025codiffmob,chu2023trajgdm}.
Real-world routes between a fixed origin and destination are diverse; they are neither purely geometric curves nor deterministic shortest paths, as route selection often depends on temporal--spatial context---time of day, surrounding land use, and perceived conditions along candidate paths.
However, the majority of existing works generate such paths by primarily using road topology features as constraints~\citep{wei2024diffrntraj,zhu2024controltraj,rao2025seed,guo2025cardiff}.
Specifically, road connectivity, intersection bottlenecks, and the hierarchical attributes of roads (e.g., highways, main roads, and residential roads) are the most common elements serving as constraints for route generation.

While such methods can generate routes that match aggregate traffic flow or cover specific regions, they fail to capture how \emph{the same origin--destination pair yields different route choices under different contexts}---a phenomenon we term \textbf{corridor diversification}.
In practice, travelers do not randomly diversify; they systematically shift toward alternative corridors in response to temporal--spatial cues (e.g., rush-hour congestion, nighttime safety concerns, or proximity to certain land uses), which results in the emergence of distinct corridors with context-dependent flow distributions.
A few works have incorporated elements related to the physical and social urban environment into route generation~\citep{cao2025hoser,su2024uen}.
However, these methods treat context as auxiliary features rather than as the mechanism that governs corridor choice, and thus fail to estimate how corridor distributions shift under different conditions.

To address these challenges, we propose \textbf{CascadeTraj}, a hierarchical route generation framework that treats temporal--spatial semantics as the conditioning signal for corridor choice, not merely a regularization feature.
The key insight is that corridor selection is a \emph{discrete decision} that should respond to context, while within-corridor trajectory execution is a continuous generation problem.
We operationalize this with two stages:
(1) the \textbf{decision stage}, which selects a corridor on the road network graph conditioned on semantic context and time;
(2) the \textbf{execution stage}, which generates a realistic trajectory within the selected corridor via residual diffusion.
As diagnostic evidence, we show that map-free continuous generation fails at corridor-level convergence---waypoints remain far from ground-truth corridors even with soft road priors (Sec.~\ref{sec:results})---motivating the discrete corridor formulation and demonstrating that semantic conditioning shifts corridor distributions under fixed OD.

The remainder of the paper develops this framing.
Section~\ref{sec:related-work} positions our work against recent trajectory foundation models~\citep{zhu2025unitraj,lin2024trajfm}, map-conditioned generators~\citep{tao2025map2traj,wang2025gtg}, and road-constrained diffusion---arguing that reconstruction fidelity is not the same as behavioral mechanism.
Section~\ref{sec:method} introduces CascadeTraj as a decision--execution framework where temporal--spatial semantics condition corridor choice.
Sections~\ref{sec:experiments}--\ref{sec:results} establish the diagnostic evidence: waypoint audits and case visualizations demonstrate the measure-zero failure mode of map-free generation, motivating the discrete corridor formulation.
